\section{Governance and Regulation}

\begin{frame}{The EU AI Act: Risk Tiers}

    {\small Good practice becomes a \textbf{requirement} once a system falls into a regulated tier.
    Orientation, not legal advice.}

    \vspace{0.2em}
    \centering
    \begin{tikzpicture}
        % Pyramid
        \fill[raiRed!70]   (0,2.6)  -- (-0.8,1.5)  -- (0.8,1.5)  -- cycle;
        \fill[raiAmber!60] (-0.8,1.5) -- (0.8,1.5) -- (1.6,0.4) -- (-1.6,0.4) -- cycle;
        \fill[raiBlue!45]  (-1.6,0.4) -- (1.6,0.4) -- (2.4,-0.7) -- (-2.4,-0.7) -- cycle;
        \fill[raiGreen!35] (-2.4,-0.7) -- (2.4,-0.7) -- (3.2,-1.8) -- (-3.2,-1.8) -- cycle;
        \draw[thick] (0,2.6) -- (3.2,-1.8) -- (-3.2,-1.8) -- cycle;

        \node[font=\tiny\bfseries, white] at (0,1.95) {Unacceptable};
        \node[font=\tiny\bfseries] at (0,0.95)  {High risk};
        \node[font=\tiny\bfseries] at (0,-0.15) {Limited risk};
        \node[font=\tiny\bfseries] at (0,-1.25) {Minimal risk};

        % Annotations
        \node[right, font=\scriptsize, text width=6.4cm, align=left] at (3.5,2.0)
            {\textbf{Prohibited (Art.~5).} Social scoring; untargeted scraping of facial images to
             build recognition databases; emotion inference at work or school; manipulative or
             exploitative systems.};
        \node[right, font=\scriptsize, text width=6.4cm, align=left] at (3.5,0.75)
            {\textbf{Permitted with obligations.} Employment, credit, education, essential
             services, law enforcement, biometrics, and safety components of regulated products
             (Annexes I and III).};
        \node[right, font=\scriptsize, text width=6.4cm, align=left] at (3.5,-0.4)
            {\textbf{Transparency duties (Art.~50).} Tell people they are interacting with an AI
             system; label synthetic media.};
        \node[right, font=\scriptsize, text width=6.4cm, align=left] at (3.5,-1.5)
            {\textbf{No specific obligations.} Spam filters, low-stakes recommendation, most
             internal tooling --- where most course projects sit.};
    \end{tikzpicture}

    \vspace{0.4em}
    {\scriptsize Source: Regulation (EU) 2024/1689 (Artificial Intelligence Act),
    \href{https://artificialintelligenceact.eu/}{artificialintelligenceact.eu}.}
\end{frame}

\begin{frame}[allowframebreaks]{The EU AI Act: What a High-Risk Provider Must Produce}

    \begin{columns}[T]
        \begin{column}{0.52\textwidth}
            \begin{itemize}\scriptsize
                \setlength\itemsep{0.1em}
                \item A \textbf{risk management system} maintained across the lifecycle.
                \item \textbf{Data governance}: relevant, representative, and examined for bias.
                \item \textbf{Technical documentation} and automatic \textbf{logging}.
                \item \textbf{Instructions for use} that state accuracy, limitations, and intended
                      conditions.
                \item Effective \textbf{human oversight} by design.
                \item Appropriate \textbf{accuracy, robustness and cybersecurity}.
                \item \textbf{Post-market monitoring} and \textbf{serious-incident reporting}.
            \end{itemize}
        \end{column}
        \begin{column}{0.46\textwidth}
            \begin{tcolorbox}[colback=raiGreen!5, colframe=raiGreen!60,
                              title=Recognise anything?]
                \scriptsize
                Disaggregated bias evaluation, documented intended use, stated limitations,
                logged decisions, monitoring.

                \vspace{0.12em}
                \textbf{That is the model card and the audit from the last two sections}, made
                mandatory. The work does not change; the audience does.
            \end{tcolorbox}

            \vspace{0.18em}
            \begin{tcolorbox}[colback=raiRed!5, colframe=raiRed!70, title=Penalties (Article 99)]
                \scriptsize
                Up to \textbf{EUR 35M or 7\%} of worldwide annual turnover for prohibited
                practices; \textbf{EUR 15M or 3\%} for most other operator obligations;
                \textbf{EUR 7.5M or 1\%} for supplying incorrect information (whichever is higher,
                with lower caps for SMEs).
            \end{tcolorbox}
        \end{column}
    \end{columns}

    \vspace{0.15em}
    {\scriptsize Regulation (EU) 2024/1689, Chapter III and Article 99. Entered into force 1 August
    2024; prohibitions and AI-literacy duties applicable from 2 February 2025; general-purpose AI
    model obligations from 2 August 2025. The 2026 ``Digital Omnibus'' amendment deferred the
    high-risk obligations (Annex III to 2 December 2027; Annex I to 2 August 2028) --- check the
    consolidated text, as these dates have already moved once.}
\end{frame}

\begin{frame}{NIST AI Risk Management Framework}

    \begin{columns}[c]
        \begin{column}{0.44\textwidth}
            \centering
            \begin{tikzpicture}[font=\scriptsize]
                \node[circle, draw, thick, fill=raiAmber!25, minimum size=1.6cm, align=center]
                    (gov) at (0,0) {\textbf{GOVERN}};
                \node[rectangle, draw, rounded corners, fill=raiBlue!15, minimum width=1.7cm,
                      minimum height=0.7cm] (map) at (0,1.8) {\textbf{MAP}};
                \node[rectangle, draw, rounded corners, fill=raiBlue!15, minimum width=1.7cm,
                      minimum height=0.7cm] (meas) at (2.0,-1.1) {\textbf{MEASURE}};
                \node[rectangle, draw, rounded corners, fill=raiBlue!15, minimum width=1.7cm,
                      minimum height=0.7cm] (man) at (-2.0,-1.1) {\textbf{MANAGE}};
                \draw[-{Latex[length=2mm]}, thick] (map) -- (meas);
                \draw[-{Latex[length=2mm]}, thick] (meas) -- (man);
                \draw[-{Latex[length=2mm]}, thick] (man) -- (map);
                \draw[gray, dashed] (gov) -- (map);
                \draw[gray, dashed] (gov) -- (meas);
                \draw[gray, dashed] (gov) -- (man);
            \end{tikzpicture}
        \end{column}
        \begin{column}{0.54\textwidth}
            \begin{itemize}\scriptsize
                \setlength\itemsep{0.4em}
                \item \textbf{Govern} --- policies, roles, accountability. Cuts across the other
                      three rather than sitting in sequence with them.
                \item \textbf{Map} --- establish the context: who is affected, what could go wrong,
                      what is the system actually for.
                \item \textbf{Measure} --- test it: metrics, disaggregated evaluation, red-teaming,
                      documented uncertainty.
                \item \textbf{Manage} --- prioritise, mitigate, monitor, and decide when to stop.
            \end{itemize}
            \vspace{0.5em}
            {\scriptsize \textbf{Voluntary}, sector-neutral, and widely used as the internal
            scaffolding for an AI governance programme --- it tells you \emph{how} to organise risk
            management, where the AI Act tells you \emph{what} you must do. ISO/IEC 42001:2023 plays
            a similar role as a certifiable management-system standard. A Generative AI Profile
            (NIST-AI-600-1, 2024) adapts the framework to generative systems.}
        \end{column}
    \end{columns}

    \vspace{0.6em}
    {\scriptsize Source: NIST, ``Artificial Intelligence Risk Management Framework (AI RMF 1.0),''
    NIST AI 100-1, January 2023,
    \href{https://www.nist.gov/itl/ai-risk-management-framework}{nist.gov}.}
\end{frame}

\begin{frame}{Organisational Practice}

    \begin{columns}[T]
        \begin{column}{0.49\textwidth}
            \textbf{Review before release}
            \begin{itemize}\small
                \setlength\itemsep{0.3em}
                \item A named reviewer who is \textbf{not} the model's author signs off the card
                      and the disaggregated results.
                \item Scale the review to the stakes: a recommender for internal search does not
                      need a committee; a credit model does.
                \item Record the decision \emph{and the trade-off accepted}, so it can be revisited.
            \end{itemize}

            \vspace{0.7em}
            \textbf{Red-teaming}
            \begin{itemize}\small
                \setlength\itemsep{0.3em}
                \item Deliberately try to make the system fail for a specific group or input type,
                      before users do.
                \item Give the red team the failure taxonomy from this lecture as a starting
                      checklist.
            \end{itemize}
        \end{column}
        \begin{column}{0.49\textwidth}
            \textbf{After release}
            \begin{itemize}\small
                \setlength\itemsep{0.3em}
                \item \textbf{Incident reporting.} A defined path for ``the model did something
                      harmful,'' with an owner and a deadline --- not a support ticket queue.
                \item \textbf{Contestability.} A person affected by a decision can ask for it to be
                      reviewed by a human, and the request is logged.
                \item \textbf{Scheduled re-audit.} Put the date in the model card and treat it like
                      a certificate expiry.
                \item \textbf{A kill switch.} Know in advance how you turn the model off and what
                      the fallback is.
            \end{itemize}
        \end{column}
    \end{columns}

    \vspace{0.8em}
    \begin{block}{}
        \centering
        \small The common failure is not the absence of a policy. It is a policy with no owner and
        no date.
    \end{block}
\end{frame}

\begin{frame}{Agents and LLMs: Covered in the NLP Course}

    Everything in this lecture applies to a model that maps an input to a prediction. Systems that
    \textbf{take actions} --- call tools, browse, write files, spend money --- add a distinct threat
    model that this deck deliberately does not duplicate.

    \vspace{0.9em}

    \begin{columns}[T]
        \begin{column}{0.49\textwidth}
            \begin{tcolorbox}[colback=gray!8, colframe=gray!60,
                              title={NLP course, Lecture 12: AI Safety for Agents}]
                \scriptsize
                \begin{itemize}
                    \item Prompt injection, direct and indirect.
                    \item Excessive agency and tool misuse.
                    \item Sandboxing, least privilege, and human-in-the-loop gates.
                    \item Agent-specific benchmarks and red-teaming.
                    \item OWASP Top 10 for LLM Applications.
                \end{itemize}
            \end{tcolorbox}
        \end{column}
        \begin{column}{0.49\textwidth}
            \begin{tcolorbox}[colback=raiBlue!5, colframe=raiBlue!60, title=What carries over]
                \scriptsize
                \begin{itemize}
                    \item Disaggregated evaluation --- an agent can fail more often for some users
                          or some languages.
                    \item Documentation of intended and out-of-scope use.
                    \item The same governance frameworks: the EU AI Act and NIST AI RMF cover both,
                          and the AI Act adds obligations for general-purpose AI models.
                \end{itemize}
            \end{tcolorbox}
        \end{column}
    \end{columns}

    \vspace{0.7em}
    \centering
    \textit{If you are shipping an agent, take both lectures. Neither is a substitute for the other.}
\end{frame}
